\phantomsection\label{sec:projects}
\section*{项目经历}

\subsection*{项目负责人}
以下项目由我担任项目负责人或主要开发者，负责整体研究方向、系统设计与实现。

\pubentry{中央高校基本科研业务费, 2021 年至 2023 年}{}
{基于轨迹表示学习的用户出行目的地与出行时间预测}
{}
{利用表示学习技术，将原始轨迹特征转换为高层信息，提升出行时间和目的地预测等下游任务的性能。}

\pubentry{个人兴趣项目, 2023 年至今}{\href{https://www.overleafcopilot.com/}{主页} \,/\, \href{https://chromewebstore.google.com/detail/overleaf-copilot/eoadabdpninlhkkbhngoddfjianhlghb}{安装}}
{OverleafCopilot: Empowering Academic Writing in Overleaf with Large Language Models}
{}
{一款将大语言模型集成到学术写作平台 Overleaf 中的浏览器扩展。}

\pubentry{个人兴趣项目, 2023 年至今}{\href{https://www.promptgenius.site/}{网站} \,/\, \href{https://github.com/wenhaomin/ChatGPT-PromptGenius}{代码}}
{PromptGenius: All-purpose Prompts for LLMs}
{}
{一个提供多种类别的大语言模型提示词的平台，支持不同的大语言模型任务。}

\subsection*{参与项目}
在以下项目中，我作为主要参与人，负责模型开发、数据分析或实验设计等具体工作。

\pubentry{Villum Foundation, 2024 年至 2026 年}{}
{基于扩散概率模型的材料逆向设计}
{}
{开发扩散概率模型以理解材料的化学成分、结构和性能之间的关系，进而对具有特定性能的新型材料开展逆向设计。}

\pubentry{中国国家自然科学基金, 2023 年至 2025 年}{}
{面向交通预测的时空轨迹数据预训练表示学习方法}
{}
{研究时空轨迹数据的预训练表示学习方法，通过建模多种特征提升机器学习方法的交通预测性能。}

\pubentry{中国国家自然科学基金, 2024 年至 2026 年}{}
{面向稀疏问题的时空轨迹生成与表示学习方法}
{}
{研究如何生成高质量时空轨迹数据及相应表示，以解决实际应用中的数据稀疏问题。}
